\subsection{純虚境界の二族の全値排除}
\label{subsec:cubic-as4-boundary-J}
\cite{AS4lowpinned}の純虚境界FI2から次の具体表示を参照する。
$\lambda=i$, $\eta=\pm1$ として
\[
\begin{gathered}
 x_0=x_1=x_2=x_5=1,\qquad
 x_{10}=x_{16}=\eta(1-i)/8,\quad
 x_{15}=-x_{11}=\eta(1+i)/8,\\
 x_{12}=h,\quad x_{17}=ih,\quad
 x_{13}=-\eta(1+i)h/4,\quad x_{18}=\eta(1-i)h/4,\\
 x_{14}=h^2+i/32,\quad x_{19}=i(h^2-1/32),\quad x_{20}=-h^2+i/32
\end{gathered}
\]
とし、他の $x_j$ は零とする。この代数を $J^\eta_h$ と書く。

\begin{proposition}\label{prop:cubic-boundary-J}
両符号の $J^\eta_h$ は、全 $h\ne0$ でH3包絡にならない。
$h=0$ の表示も14641型AS4ではなく、H3から排除される。
\end{proposition}
\begin{proof}
両符号を別々に $\mathbb Q(i)(h)$ 上で計算する。
最高tensor $R A R^{\mathsf T}$ は全 $h$ で非零反巡回なので、AS4実現では $\mu=I$ である。
三次・四次の全関係moduleから構成した商座標は、分母が $h$ の冪のみであり、
全関係を消し自由列上でidentityとなる。
従って $h\ne0$ のAS4実現では真の20,35次元基底である。
中心三次行列の20次minorはそれぞれ
\[
 \eta=+1:\quad \frac{i h^9}{2^{11}},\qquad
 \eta=-1:\quad \frac{(1-i)h^{10}}{2^7}
\]
であり、全非零 $h$ で階数20を得る。
従って中心三次が零となり、H3を排除する。

$h=0$ には上の商座標を適用しない。
元の六関係の五次全語行列を $\mathbb Q(i)$ 上で別に構成すると、
両符号でその階数は966であり、
\[
 \dim(J^\eta_0)_5=4^5-966=58>56=\binom83.
\]
14641型AS4のHilbert級数に反するので、こちらもH3にならない。
四次までのHilbert値が一致し得ることを、AS4十分性として使っていない。
\end{proof}

二符号を直接検算し、分母支持と端点の五次行列も確認する再現コードは
\texttt{verify\_\allowbreak{}cubic\_\allowbreak{}as4\_\allowbreak{}four\_\allowbreak{}simple\_\allowbreak{}boundary\_\allowbreak{}J.py}
である。FI2表示のみを参照し、全chart/CY被覆・AS4存在十分性・全同型は輸入しない。
